\chapter{Introduction}
\label{cha:intro}

\section{Background and Motivation}
Bayesian inference provides a rigorous framework for updating beliefs 
about model parameters in the light of observed data, yet the uncertainty 
it quantifies is not immune to uncertainty itself. This is the problem of 
second-order uncertainty, which arises because Bayesian inference carries a 
foundational assumption that the prior distribution is known and correctly 
specified. When multiple priors are equally plausible, committing to a single 
one is epistemically insufficient and can lead to overconfident inferences. 
The framework is precise about uncertainty over parameters but it fails to 
capture the uncertainty about the choice of the prior which precedes inference 
at the level of modelling itself. This problem has been recognised in the 
philosophy of probability since at least the 1960s, formalised under the 
study of imprecise probabilities and credal sets. Imprecise probability theory 
replaces the single prior distribution with a credal set, a convex set of 
distributions representing the range of plausible beliefs, and asks what 
conclusions are robust across the entire set rather than what is optimal 
under a single prior \citep{walley1991}\citep{wheeler2022}. This is a philosophically 
distinct response to the problem of second-order uncertainty. Rather than 
resolving prior uncertainty by collapsing it to a point, the framework carries 
that uncertainty forward explicitly into the inferential conclusions. The 
computational challenges of working with sets of distributions are formidable, 
however, and the focus on scalability in modern machine learning has led to 
the problem of second-order uncertainty being largely neglected.

Prior sensitivity is a routine challenge in Bayesian inference. When data 
is scarce, the likelihood is not sufficiently informative to overwhelm the 
influence of the prior and the posterior remains sensitive to the initial 
modelling choices. In settings such as healthcare, ecological modelling, or 
scientific parameter estimation, data is expensive to obtain or inherently 
limited; two practitioners using equally defensible priors can arrive at 
substantially different posteriors and therefore different conclusions. In a 
clinical trial setting, for instance, a sceptical prior and an optimistic 
prior over treatment efficacy can produce posteriors that disagree 
substantially about whether a drug should be recommended, not because the 
data is ambiguous, but because the prior is doing significant inferential 
work that is never formally acknowledged \citep{berger1994}. This sensitivity 
is not confined to small data or traditional statistical settings. Recent 
work has shown that prior choice remains consequential in deep Bayesian 
models; standard priors used in Bayesian neural networks are poorly 
calibrated and performance is sensitive to prior specification even at 
scale \citep{fortuin2021}, and cold posteriors, which correspond to 
implicitly rescaled priors, have been shown to outperform standard 
posteriors in deep learning settings \citep{wenzel2020}. Prior 
uncertainty is therefore not a historical concern that modern methods 
have resolved but an active and open problem in probabilistic machine 
learning.

Variational Inference (VI) has emerged as a powerful tool for Bayesian 
computation, offering a scalable alternative to Markov Chain Monte Carlo 
(MCMC) methods. However, standard VI inherits the same foundational 
limitation: it approximates a single posterior under a fixed prior, leaving 
second-order uncertainty unaddressed. VI is nevertheless a natural framework 
for operationalising credal sets at scale: rather than running a separate 
inference procedure for each candidate prior independently, one can train a 
family of variational posteriors jointly and combine them into a credal 
mixture, directly representing the uncertainty that a single posterior cannot 
capture. The computational challenge is making this tractable; training 
expressive posterior approximations across a set of priors simultaneously 
demands a flexible and scalable variational family.

Recent advances in normalising flows and their continuous-time generalisation 
through neural ordinary differential equations \citep{chen2018} have made 
expressive posterior approximation tractable. Flow matching \citep{lipman2022} represents a further step, replacing the computationally expensive 
log-determinant evaluations required by classical normalising flows with a 
direct regression objective over a learned velocity field. Crucially, the 
velocity field can be conditioned on auxiliary information, in this case a 
discrete prior index, at minimal additional cost, allowing a single shared 
model to approximate posteriors across all candidate priors simultaneously. 
This makes flow matching particularly well suited to the credal mixture 
setting and addresses the computational challenge directly.

This thesis investigates whether credal mixtures of flow-based posteriors 
can provide robust inference under prior uncertainty. The approach is grounded 
in a classical result from information theory: the KL divergence is convex 
in its second argument, which guarantees in principle that a mixture over a 
set of posteriors performs at least as well as the best individual component 
\citep{cover2006}. Whether this guarantee remains empirically visible 
under deep variational approximation, and what architectural and methodological 
conditions are required for it to hold, is the central question this work 
addresses.

\section{Existing Approaches to Prior Robustness}

Several approaches have been proposed to address the problem of prior 
sensitivity in Bayesian inference and to operationalise the concept of 
robustness to prior uncertainty. Each offers a different perspective on 
the problem and a different set of trade-offs, but none provides a 
principled treatment of second-order uncertainty in the sense formalised 
by imprecise probability theory.

\subsection{Epsilon-Contamination and Prior Classes}

The original formal treatment of prior robustness models the prior as 
belonging to an epsilon-contamination class, a convex set of distributions 
formed by mixing a baseline prior with an arbitrary contaminating 
distribution within some neighbourhood (Berger, 1984; Insua and Ruggeri, 
2000). Posterior bounds are then derived that hold for any prior in the 
class, providing a formal framework for analysing the sensitivity of 
inferences to perturbations of the prior. This is the intellectual 
precursor to the credal set approach and establishes the core idea that 
working with a class of priors rather than a single one is both principled 
and tractable. The limitation is that the prior class is defined 
analytically and the resulting posterior bounds are derived through 
closed-form calculations, which restricts applicability to simple models 
and low-dimensional settings. This approach does not scale to modern 
machine learning problems with high-dimensional parameter spaces and 
complex likelihoods.

\subsection{Gamma-Minimax Decision Theory}

Rather than committing to one prior, the gamma-minimax framework minimises 
the maximum expected loss over a class of priors, producing decisions that 
are robust to the worst-case prior within the class \citep{berger1985}. This 
is philosophically aligned with the imprecise probability approach in that 
it acknowledges a set of plausible priors rather than a single one. 
However, the minimax framing collapses the set at decision time: the 
output is a single decision that hedges against the worst case, rather 
than a set of posteriors that carries prior uncertainty forward. The 
second-order uncertainty is acknowledged in the decision procedure but 
is not propagated through inference in a way that allows downstream 
conclusions to reflect it.

\subsection{Bayesian Model Averaging}

Bayesian model averaging (BMA) combines a discrete set of candidate 
models or priors by weighting each according to its marginal likelihood 
\citep{hoeting1999}. This is the closest classical precursor to the 
credal mixture approach and is important to engage with carefully. The 
key distinction is that BMA uses marginal likelihood weights, which are 
themselves computed under each model's prior and collapsed into a single 
weighted predictive distribution. The weights are point estimates and the 
resulting average is a single distribution rather than a credal set. There 
is no formal commitment to representing the uncertainty across the set; 
BMA resolves prior uncertainty by weighting rather than carrying it forward. 
Furthermore, marginal likelihood weights can be dominated by a single model 
when one prior fits the data substantially better, which may not reflect 
genuine robustness across the set.

\subsection{Stacking of Predictive Distributions}

Stacking fits mixture weights over a set of candidate posteriors by 
maximising held-out log predictive scores, making it the most directly 
related existing method to the approach developed in this thesis \citep{yao2018}. Operationally, stacking shares the combination mechanism: 
a set of posteriors is trained and weights are selected using a 
likelihood-consistent objective. However, stacking treats the candidate 
models as fixed and selects weights purely for predictive performance, 
without any theoretical grounding in imprecise probability theory or the 
KL convexity guarantee. It is a pragmatic ensemble method rather than 
a principled treatment of second-order uncertainty. The mixture weights 
in stacking are justified by cross-validated predictive accuracy alone; 
in the credal mixture framework developed here, the combination is 
grounded in the convexity of KL divergence, which provides a theoretical 
guarantee that the mixture performs at least as well as the best 
individual component.

\subsection{Generalised Bayesian Inference}

Generalised Bayesian inference replaces the likelihood in Bayes' rule 
with a general loss function, producing a family of posteriors indexed 
by the loss rather than a single posterior \citep{bissiri2016}. 
Knoblauch et al. \citep{knoblauch2022} extend this framework and explicitly discuss 
robustness to both prior and likelihood misspecification as a unified 
problem, making this the most theoretically sophisticated modern approach 
to Bayesian robustness. The framework implicitly defines a set of 
plausible posteriors under different modelling assumptions, which is 
conceptually adjacent to the credal set idea. However, it does not 
formally invoke imprecise probability theory; it does not treat the 
resulting set of posteriors as a credal set or ask what can be concluded 
robustly across it. It produces alternative posteriors rather than 
representing and propagating genuine second-order uncertainty, and it 
does not provide a mechanism for combining posteriors across a set 
of priors in a principled way.

\subsection{Robust Divergence-Based Variational Inference}

Robust divergence-based variational inference replaces the KL divergence 
in the VI objective with a more robust alternative such as the 
beta-divergence or gamma-divergence, achieving resistance to outliers 
and likelihood misspecification \citep{futami2018, knoblauch2019}. While superficially related to the goal of robust inference, 
this approach addresses a fundamentally different problem. It produces 
a single posterior under a single prior, changing the objective function 
to be less sensitive to data contamination but leaving the question of 
which prior to use entirely unaddressed. A practitioner using these 
methods still commits to one prior at the outset and obtains one 
posterior at the end; the second-order uncertainty identified in 
Section 1.1 is not represented, acknowledged, or propagated.

\bigskip

None of these approaches simultaneously satisfies three properties 
that the present work aims to achieve: formal grounding in imprecise 
probability theory, expressive deep variational posterior approximation, 
and likelihood-consistent combination of posteriors under a credal set 
of priors. The closest in terms of combination mechanism is stacking, 
which lacks the theoretical grounding. The closest theoretically is 
generalised Bayesian inference, which lacks both the credal set framing 
and the combination mechanism. Imprecise probability theory is 
philosophically distinct from all of the above: by treating the prior 
as a credal set from the outset, it propagates uncertainty at the 
modelling level through to the posterior, asking what can be concluded 
robustly across all plausible priors rather than optimising under any 
single one \citep{walley1991, wheeler2022}. This thesis operationalises 
that commitment computationally for the first time in the setting of 
flow-based variational inference.

\section{Problem Statement}

Having established that existing approaches fall short of a principled 
treatment of second-order uncertainty, we can now state the specific 
problem this thesis investigates more precisely. The central challenge 
is the gap between the theoretical promise of credal mixture inference 
and its empirical behaviour under deep variational approximation. While 
the convexity of the KL divergence guarantees that a mixture over a 
set of posteriors should perform at least as well as the best individual 
component, it is not clear whether this guarantee remains empirically 
visible when using deep neural networks to approximate the posteriors. 
These approximations are learned under a finite data regime and the 
mixture weights must be estimated from held-out data. These sources of 
approximation and estimation error may undermine the theoretical 
guarantee, and it is an open question under what conditions the robust 
performance of credal mixtures can be realised in practice.

More concretely, the problem has two interacting components. The first 
is architectural: if each posterior approximation is trained independently 
under a different prior, there is no guarantee that the resulting 
approximations are comparable in quality or that their mixture is 
well-defined as a density. The second is statistical: even with a 
coherent family of approximations, estimating mixture weights using 
proxy objectives such as maximum mean discrepancy or adversarial 
divergence bounds introduces estimation noise that can cause the mixture 
to underperform its best component. Both failure modes were observed in 
preliminary experiments and motivate the specific design choices 
investigated in this thesis.

\section{Research Objectives}

This thesis pursues four concrete objectives in response to the 
challenges identified above.

\begin{enumerate}
    \item \textbf{Shared-backbone architecture.} Design a flow matching 
    architecture conditioned on a discrete prior index, ensuring that 
    all posterior approximations are trained within a common optimisation 
    framework with comparable capacity across priors.

    \item \textbf{Likelihood-consistent weight selection.} Investigate 
    held-out negative log likelihood as a principled alternative to proxy 
    divergence objectives such as maximum mean discrepancy and adversarial 
    KL bounds, which were found in preliminary experiments to introduce 
    estimation noise sufficient to mask the theoretical guarantee. Unlike 
    proxy objectives, held-out NLL directly approximates the KL objective 
    underpinning the convexity guarantee, ensuring alignment between the 
    selection criterion and the theoretical motivation.

    \item \textbf{Empirical evaluation of the convexity guarantee.} 
    Evaluate whether the resulting credal mixtures recover the KL 
    convexity guarantee empirically, using paired per-seed comparisons 
    to control for estimation variance across random initialisations.

    \item \textbf{Characterisation of failure conditions.} Identify the 
    architectural and statistical conditions under which the theoretical 
    guarantee is and is not empirically visible, contributing to a broader 
    understanding of when classical information-theoretic arguments remain 
    reliable under deep variational approximation.
\end{enumerate}

Taken together, these objectives aim to bridge the gap between the 
theoretical commitments of imprecise probability theory and their 
empirical behaviour under deep variational approximation, contributing 
to a broader understanding of when information-theoretic guarantees 
survive the transition from exact to approximate inference.

\section{Report Structure}

The remainder of this thesis is organised as follows.

Chapter~2 presents the theoretical foundations 
underpinning the approach. It introduces Bayesian 
inference and variational inference, formalises 
imprecise probability theory and credal sets, and 
establishes the KL convexity guarantee that motivates 
the credal mixture framework. It then develops the 
normalising flow and flow matching machinery, and 
proves the approximate mixture bound that characterises 
the conditions under which the guarantee survives 
flow-based approximation and finite held-out weight 
selection.

Chapter~3 describes the methodology. It formalises 
the three research hypotheses and the two-regime 
experimental design, introduces the datasets and 
credal sets used for each problem, and describes the 
shared-backbone \texttt{VelocityResNet} architecture 
and its prior conditioning mechanism. The two-phase 
training procedure, the eight credal weight 
optimisation algorithms compared, and the full 
evaluation protocol are also defined here.

Chapter~4 presents the results and discussion. 
The primary evaluation on German Credit in the 
low-data regime confirms H1 and H3, demonstrating 
that the credal mixture achieves near-oracle 
worst-case regret and that EM is the most consistent 
weight optimisation algorithm. The full-data null 
control confirms H2. An architectural ablation 
validates the shared-backbone design, and 
supplementary analyses examine prior sensitivity, 
weight stability, leave-one-prior-out behaviour, 
calibration, and optimiser convergence.Section~\ref{sec:fm_vs_cnf} presents a 
speed comparison between the flow matching 
training objective and a CNF baseline using 
the same architecture, validating the 
computational motivation of Section~2.5.

Chapter~5 concludes the thesis with a summary of 
contributions, an acknowledgement of limitations, 
and directions for future work including continuous 
credal sets, theoretical bounds on the approximation 
gap, and extension to Bayesian neural networks.